\icode{E-ME-MU-L3-261019-001}
如图所示，粗细均匀的光滑直杆竖直固定，小球\(A\)套在杆上，球孔的内径比杆直径略大，绕过定滑轮的轻绳一端连接在小球\(A\)上，另一端连着质量为\(m\)的物块\(B\)，开始时，连接\(A\)的轻绳与水平方向的夹角为\qty{37}{\degree}，定滑轮到小球\(A\)间的轻绳长度为\(L\)，杆上的\(O\)点与定滑轮等高，杆上\(C\)点与定滑轮的连线与水平方向的夹角为\qty{37}{\degree}，\(\sin\)\qty{37}{\degree}=0.6，\(\cos\)\qty{37}{\degree}=0.8，由静止释放小球\(A\)后，小球\(A\)向上运动，不计小球和定滑轮的大小，重力加速度为\(g\)，则下列判断正确的是\paren[]

\begin{tikzpicture}[font=\scriptsize]
	\begin{scope}[rotate=26.5]
		\draw (0,0) circle(.1);
		\foreach \x in {-.6,-.4,...,.6}
			\draw[gray] (\x,.8)--++(.2,-.2);
		\draw[thick] (-.6,.6)--(.6,.6) (0,0)--(0,.6);
		\fill (0,0) circle(.05);
	\end{scope}
	\draw[magenta] (53:.1)--++(-37:3)coordinate(A) (-.1,0)--(-.1,-1.2);
	\coordinate (e) at($(A)+(-.05,-.5)$);
	\coordinate (f) at($(A)+(.05,4)$);
	\coordinate (g) at(A|-0,.1);
	\coordinate (C) at($(A)!2!(g)$);
	\draw[dashed,cyan] (0,.1)--(g)node[right,black]{\(O\)};
	\draw[fill=orange!40] (e) rectangle(f) (-.3,-1.2) rectangle(.1,-1.6)node[midway]{\(B\)};
	\shade[ball color=gray] (A)node[right]{\(A\)} circle(.12);
	\fill (C)node[right]{\(C\)} circle(.05) (g) circle(.05);
	\node at(-.3,-1.4)[left]{\(m\)};
\end{tikzpicture}
\begin{choices}
	\item 若小球恰好能到达\(O\)点，则小球\(A\)的质量为\(\frac{1}{3}m\)
	\item 若小球恰好能到达\(O\)点且到达\(O\)点时受到竖直向上的恒力，之后小球恰好能到达到\(C\)点，则该恒力的大小为\(\frac{1}{2}mg\)
	\item 小球 A 有可能恰好运动到 C 点
	\item 若小球\(A\)的质量为\(0.5m\)，则小球\(A\)向上运动过程中，物块\(B\)先加速后减速
\end{choices}
\begin{answer}
    
\end{answer}